%% toc: P15 157 157.7
% Proofs for the source variants of the measure-theoretic foundations (Durrett, Probability: Theory and Examples, Version 5, pp. 2--3, 6--17, 53--54, 210, 214--215, 457--458).
\section{Weights, generators, extended limits and coding variants: proofs}
\label{sec:annex-p15-w02c-p15-002-source-variants}

Use the definitions of
Chapter~\ref{sec:p15-w02c-p15-002-source-variants} and of
\cref{sec:p15-measurable-probability}. Two facts about the real numbers
are used repeatedly: between any two distinct real numbers there is a
rational number, and for every real \(x\) there are rationals below and
above \(x\).

\begin{annexitem}{annex:p15-w02c-p15-002-generator-variants}{thm:p15-w02c-p15-002-generator-variants}{Open rational rays and real open rectangles generate the Borel sets}
\begin{proof}
(i) Let \(\mathcal G_1=\sigma(\{(-\infty,x):x\in\mathbb Q\})\). Each ray
\((-\infty,x)\) is open: if \(y<x\), the open ball of radius \(x-y\) about
\(y\) lies in the ray. Hence each generator is Borel and
\(\mathcal G_1\subseteq\mathcal B(\mathbb R)\). Conversely fix a real
\(a\) and, for each \(n\geq1\), choose a rational \(x_n\) with
\(a<x_n<a+1/n\). If \(y\leq a\), then \(y<x_n\) for every \(n\). If
\(y>a\), choose \(n\) with \(1/n<y-a\); then \(x_n<a+1/n<y\), so
\(y\notin(-\infty,x_n)\). Therefore
\[
 (-\infty,a]=\bigcap_{n\geq1}(-\infty,x_n)\in\mathcal G_1 .
\]
By \cref{thm:p15-borel-generators} the rays \((-\infty,a]\), \(a\) real,
generate \(\mathcal B(\mathbb R)\), so
\(\mathcal B(\mathbb R)\subseteq\mathcal G_1\).

(ii) Let \(\mathcal G_2\) be the sigma-algebra generated by the open
rectangles \(R=\prod_{j=1}^d(a_j,b_j)\) with real \(a_j<b_j\). Each such
\(R\) is open: for \(x\in R\) put
\(r=\min_j\min(x_j-a_j,\,b_j-x_j)>0\); if the Euclidean distance
\(|y-x|\) is less than \(r\), then \(|y_j-x_j|\leq|y-x|<r\) for every
coordinate, so \(a_j<y_j<b_j\) and \(y\in R\). Thus
\(\mathcal G_2\subseteq\mathcal B(\mathbb R^d)\). The rectangles with
rational endpoints are among the generators of \(\mathcal G_2\), and they
generate \(\mathcal B(\mathbb R^d)\) by \cref{thm:p15-borel-generators};
hence \(\mathcal B(\mathbb R^d)\subseteq\mathcal G_2\). The class of open
sets generates \(\mathcal B(\mathbb R^d)\) by
\cref{def:p15-borel-sets}, and it contains every such rectangle by what
was just shown.

(iii) Apply \cref{thm:p15-measurability-criterion} with the generating
class of (i), noting \(X^{-1}((-\infty,x))=\{X<x\}\), and with the
generating class of (ii).
\end{proof}
\end{annexitem}

\paragraph{Difference operators.} For \(G:\mathbb R^d\to\mathbb R\), a
coordinate \(k\), reals \(s<t\) and \(x\in\mathbb R^d\), let \(x[k\gets t]\)
be \(x\) with its \(k\)-th coordinate replaced by \(t\), and
\[
 (D^k_{s,t}G)(x)\defeq G(x[k\gets t])-G(x[k\gets s]).
\]
Each \(D^k_{s,t}\) is linear, and its result does not depend on \(x_k\).
Operators acting on different coordinates commute, since both orders give
the same four signed terms. For \(s<u<t\),
\(D^k_{s,u}+D^k_{u,t}=D^k_{s,t}\), because the middle terms cancel; by
induction, for \(s=\alpha_0<\alpha_1<\cdots<\alpha_m=t\),
\(D^k_{s,t}=\sum_{j=1}^mD^k_{\alpha_{j-1},\alpha_j}\). Expanding the
composition, for a finite rectangle \(A=\prod_k(a_k,b_k]\),
\[
 \bigl(D^1_{a_1,b_1}D^2_{a_2,b_2}\cdots D^d_{a_d,b_d}F\bigr)(x)
   =\sum_{v\in V}\operatorname{sgn}(v)F(v)=\Delta_AF
 \qquad(x\in\mathbb R^d),
\tag{\(\dagger\)}
\]
since each vertex \(v\) arises from exactly one choice of \(b_k\) or
\(a_k\) in each coordinate, with one minus sign for each \(a_k\) chosen;
by commutativity the operators may be applied in any order.

\begin{annexitem}{annex:p15-w02c-p15-002-rectangle-inclusion-exclusion}{thm:p15-w02c-p15-002-rectangle-inclusion-exclusion}{The difference \(\Delta_AF\) is the mass of the rectangle}
\begin{proof}
Write \((-\infty,x]=\prod_i(-\infty,x_i]\). Fix a finite rectangle
\(A=\prod_i(a_i,b_i]\). For \(k=0,\ldots,d\) and \(x\in\mathbb R^d\) let
\[
 R_k(x)=\prod_{i\leq k}(a_i,b_i]\times\prod_{i>k}(-\infty,x_i],\qquad
 G_k(x)=\mu(R_k(x)) .
\]
These are Borel sets, being products of intervals
(\cref{thm:p15-w02c-p15-002-generator-variants} and its proof), and
\(G_k(x)\leq F(\max(x,b))<\infty\), where the maximum is taken
coordinatewise. Here \(G_0=F\) and \(G_d(x)=\mu(A)\). For \(k\geq1\),
\(R_{k-1}(x[k\gets b_k])\) is the disjoint union of
\(R_{k-1}(x[k\gets a_k])\) and \(R_k(x)\), because \((-\infty,b_k]\) is the
disjoint union of \((-\infty,a_k]\) and \((a_k,b_k]\). By part~(i) of
\cref{thm:p15-measure-laws}, and since all values are finite,
\(G_k=D^k_{a_k,b_k}G_{k-1}\). Hence
\(\mu(A)=G_d=D^d_{a_d,b_d}\cdots D^1_{a_1,b_1}F=\Delta_AF\) by
\((\dagger)\).

(iv) follows: \(\Delta_AF=\mu(A)\geq0\). (i): if \(x\leq y\) then
\((-\infty,x]\subseteq(-\infty,y]\), so \(F(x)\leq F(y)\). (ii): let
\(\mathbf 1=(1,\ldots,1)\). The sets \((-\infty,x+\mathbf1/m]\) decrease
to \((-\infty,x]\) and have finite measure, so part~(iv) of
\cref{thm:p15-measure-laws} gives \(F(x+\mathbf1/m)\downarrow F(x)\).
Given \(\varepsilon>0\) choose \(m\) with
\(F(x+\mathbf1/m)<F(x)+\varepsilon\); every \(y\) with
\(x\leq y\leq x+\mathbf1/m\) has \(F(x)\leq F(y)<F(x)+\varepsilon\) by
(i). So \(F(y)\to F(x)\) as \(y\downarrow x\).

Now let \(\mu\) be a probability measure. For fixed \(i\) and \(x\), the
sets \((-\infty,x[i\gets-m]]\) decrease in \(m\) to the empty set, so
\(F(x[i\gets-m])\to0\) by part~(iv) of \cref{thm:p15-measure-laws}; for
real \(t\leq-m\), (i) gives \(0\leq F(x[i\gets t])\leq F(x[i\gets-m])\), so
the limit as \(t\to-\infty\) is \(0\). The sets \((-\infty,m\mathbf1]\)
increase to \(\mathbb R^d\), so \(F(m\mathbf1)\to1\) by part~(iii) of the
same theorem. If every coordinate of \(x_n\) tends to \(\infty\), then for
each \(m\) eventually \(x_n\geq m\mathbf1\), so
\(F(m\mathbf1)\leq F(x_n)\leq1\) eventually, and \(F(x_n)\to1\). This is
(iii\(^*\)). Finally (iii\(^*\)) with (i) gives (iii): if every coordinate
of \(x_n\) tends to \(-\infty\), fix any \(c\in\mathbb R\); eventually all
coordinates of \(x_n\) are at most \(c\), so
\(0\leq F(x_n)\leq F((c,\ldots,c)[1\gets x_{n,1}])\to0\).
\end{proof}
\end{annexitem}

\begin{annexitem}{annex:p15-w02c-p15-002-negative-mass}{thm:p15-w02c-p15-002-negative-mass}{The counterexample and a second one}
\begin{proof}
(i) Let \(c(t)=0\) for \(t<0\), \(c(t)=1\) for \(0\leq t<1\) and
\(c(t)=2\) for \(t\geq1\). Then \(F(x)=f(c(x_1),c(x_2))\), where
\(f(2,2)=1\), \(f(2,1)=f(1,2)=2/3\), and \(f=0\) at the other six pairs.
The function \(c\) is nondecreasing, and \(f\) is nondecreasing in each
argument: \(f(0,\cdot)\) and \(f(\cdot,0)\) are \(0\), \(f(1,\cdot)\) and
\(f(\cdot,1)\) take the values \(0,0,2/3\), and \(f(2,\cdot)\) and
\(f(\cdot,2)\) take \(0,2/3,1\). Hence (i). Each set
\(\{t:c(t)=j\}\) is an interval closed on the left and open on the right,
so if \(y_i\downarrow x_i\) then \(c(y_i)=c(x_i)\) once \(y_i\) is close
enough to \(x_i\); thus \(F(y)=F(x)\) for \(y\geq x\) near \(x\), which
gives (ii). If some \(x_i<0\), then \(c(x_i)=0\) and \(F(x)=0\), so
\(F(x[i\gets t])=0\) for \(t<0\); if both coordinates are at least \(1\),
\(F=1\). This is (iii\(^*\)), and (iii) follows from it with (i), by the
argument at the end of the proof of
\cref{thm:p15-w02c-p15-002-rectangle-inclusion-exclusion}.

For \(0<\varepsilon\leq1\), \(1-\varepsilon\in[0,1)\) and
\(-\varepsilon<0\). So \(F(1,1)=1\), \(F(1-\varepsilon,1)=F(1,1-\varepsilon)=2/3\),
\(F(1-\varepsilon,1-\varepsilon)=0\), and
\(\Delta_{(1-\varepsilon,1]^2}F=1-2/3-2/3+0=-1/3\). Likewise
\(\Delta_{(1-\varepsilon,1]\times(-\varepsilon,0]}F
 =F(1,0)-F(1-\varepsilon,0)-F(1,-\varepsilon)+F(1-\varepsilon,-\varepsilon)
 =2/3-0-0+0\), and the case \((-\varepsilon,0]\times(1-\varepsilon,1]\) is
symmetric. Since \(-1/3<0\), (iv) fails, and by
\cref{thm:p15-w02c-p15-002-rectangle-inclusion-exclusion} \(F\) is not the
distribution function of any measure.

(ii) With \(h=\mathbf1_{[0,\infty)}\), \(F_0(x)=\frac12(h(x_1)+h(x_2))\)
takes values in \(\{0,1/2,1\}\). It is nondecreasing and right continuous
because \(h\) is. If both coordinates tend to \(-\infty\), eventually both
are negative and \(F_0=0\); if both tend to \(\infty\), eventually
\(F_0=1\). So (i)--(iii) hold. The function \(x\mapsto h(x_1)\) does not
depend on \(x_2\), so \(D^2_{a_2,b_2}\) maps it to \(0\); likewise
\(D^1_{a_1,b_1}\) maps \(x\mapsto h(x_2)\) to \(0\). By linearity,
commutativity and \((\dagger)\),
\(\Delta_AF_0=\frac12\bigl(D^1D^2h(x_1)+D^2D^1h(x_2)\bigr)=0\) for every
finite rectangle, so (iv) holds. Suppose a probability measure \(\mu\) had
\(\mu(A)=\Delta_AF_0=0\) for all finite rectangles. The squares
\((-n,n]^2\) increase to \(\mathbb R^2\), so by part~(iii) of
\cref{thm:p15-measure-laws},
\(\mu(\mathbb R^2)=\lim_n\mu((-n,n]^2)=0\), contradicting
\(\mu(\mathbb R^2)=1\). Finally \(F_0(t,0)=\frac12(0+1)=1/2\) for \(t<0\).
\end{proof}
\end{annexitem}

\begin{annexitem}{annex:p15-w02c-p15-002-rd-construction}{thm:p15-w02c-p15-002-rd-construction}{Durrett's Theorem 1.1.11 with the strong tail condition}
\begin{proof}
Write \(\mathbf1=(1,\ldots,1)\), \((a,b]=\prod_i(a_i,b_i]\) for
\(a<b\) in every coordinate, and \(m(A)=\Delta_AF\) for finite rectangles
\(A\), \(m(\varnothing)=0\). By (iv), \(m\geq0\).

\emph{Step 1: additivity over finite partitions.} Call a partition of
\(A=(a,b]\) \emph{regular} if there are points
\(a_k=\alpha_{k,0}<\alpha_{k,1}<\cdots<\alpha_{k,n_k}=b_k\) for each \(k\)
such that its members are the cells
\(\prod_k(\alpha_{k,j_k-1},\alpha_{k,j_k}]\), \(1\leq j_k\leq n_k\).
These cells are pairwise disjoint and cover \(A\), since each coordinate
interval \((a_k,b_k]\) is the disjoint union of the
\((\alpha_{k,j-1},\alpha_{k,j}]\). Writing
\(D^k_{a_k,b_k}=\sum_jD^k_{\alpha_{k,j-1},\alpha_{k,j}}\) in \((\dagger)\)
and expanding by linearity gives \(m(A)=\sum_{\text{cells }C}m(C)\).
Now let \(A\) be the disjoint union of finite rectangles
\(A_1,\ldots,A_r\). For each \(k\), list in increasing order all the
\(k\)-th endpoints of \(A\) and of the \(A_j\); this gives a regular
partition of \(A\). A cell \(C\) that meets \(A_j\) lies in \(A_j\): in each
coordinate the cell interval runs between consecutive listed points and
meets the interval of \(A_j\), whose endpoints are listed points, so it lies
inside it. Each cell lies in \(A\), hence meets, and so lies in, exactly
one \(A_j\); and the cells inside \(A_j\) form a regular partition of
\(A_j\). Therefore
\(m(A)=\sum_Cm(C)=\sum_j\sum_{C\subseteq A_j}m(C)=\sum_jm(A_j)\).

\emph{Step 2: the box semialgebras.} For \(n\geq1\) let
\(B_n=(-n\mathbf1,n\mathbf1]\) and let \(\mathcal S_n\) consist of
\(\varnothing\) and the finite rectangles contained in \(B_n\). The
intersection of two members is empty or a finite rectangle in \(B_n\),
coordinatewise. For \((a,b]\in\mathcal S_n\), a point of
\(B_n\setminus(a,b]\) has a first coordinate \(k\) with
\(x_k\notin(a_k,b_k]\), so
\[
 B_n\setminus(a,b]=\bigcup_{k=1}^d\Bigl(\prod_{i<k}(a_i,b_i]
   \times\bigl((-n,a_k]\cup(b_k,n]\bigr)\times\prod_{i>k}(-n,n]\Bigr),
\]
a disjoint union of at most \(2d\) members of \(\mathcal S_n\) after
deleting empty pieces. So \(\mathcal S_n\) is a semialgebra on \(B_n\), and
by Step~1 \(m\) satisfies the finite partition condition of
\cref{thm:p15-semialgebra-extension}. That theorem gives the finitely
additive extension \(\overline m_n\) to \(\mathcal A(\mathcal S_n)\), monotone
and finitely subadditive over covers.

\emph{Step 3: countable subadditivity.} Let \((a,b]\in\mathcal S_n\) and
\((a,b]\subseteq\bigcup_{i\geq1}(a^i,b^i]\) with
\((a^i,b^i]\in\mathcal S_n\) (empty members can be dropped). Let
\(\varepsilon>0\). By \((\dagger)\) and (ii), \(m((a+\delta\mathbf1,b])\to
m((a,b])\) as \(\delta\downarrow0\), because each vertex with some
coordinates \(a_k+\delta\) decreases coordinatewise to the corresponding
vertex of \((a,b]\); choose \(0<\delta<\min_k(b_k-a_k)\) with
\(m((a+\delta\mathbf1,b])>m((a,b])-\varepsilon\). In the same way choose
\(0<\eta_i<1\) with
\(m((a^i,b^i+\eta_i\mathbf1])<m((a^i,b^i])+\varepsilon2^{-i}\). The
closed rectangle \([a+\delta\mathbf1,b]\) is compact
(\cref{thm:p15-cylinder-compact-approximation}) and lies in
\((a,b]\subseteq\bigcup_i(a^i,b^i+\eta_i\mathbf1)\), a union of open
rectangles (open by the proof of
\cref{thm:p15-w02c-p15-002-generator-variants}). Finitely many,
\(i\leq J\), cover it, so
\((a+\delta\mathbf1,b]\subseteq\bigcup_{i\leq J}(a^i,b^i+\eta_i\mathbf1]\).
All these rectangles lie in \(B_{n+1}\), so the finite subadditivity of
\(\overline m_{n+1}\) gives
\[
 m((a,b])-\varepsilon<m((a+\delta\mathbf1,b])
 \leq\sum_{i\leq J}m((a^i,b^i+\eta_i\mathbf1])
 \leq\sum_{i\geq1}m((a^i,b^i])+\varepsilon .
\]
As \(\varepsilon\) is arbitrary,
\(m((a,b])\leq\sum_im((a^i,b^i])\). By
\cref{thm:p15-semialgebra-extension}, \(\overline m_n\) is a premeasure, and
it is finite since \(\overline m_n(B_n)=m(B_n)<\infty\). By
\cref{thm:p15-premeasure-extension} it has a unique extension to a measure
\(\mu_n\) on \(\sigma(\mathcal S_n)\).

\emph{Step 4: the sigma-algebra.} Members of \(\mathcal S_n\) are Borel
subsets of \(B_n\), so \(\sigma(\mathcal S_n)\) is contained in the
collection \(\mathcal B(B_n)\) of Borel subsets of \(B_n\). Conversely
\(\mathcal L=\{E\in\mathcal B(\mathbb R^d):E\cap B_n\in\sigma(\mathcal S_n)\}\)
is a sigma-algebra on \(\mathbb R^d\) (complements: \(E^c\cap B_n=B_n\setminus
(E\cap B_n)\); unions: termwise), and it contains every finite rectangle
\(R\), since \(R\cap B_n\) is empty or in \(\mathcal S_n\). Each open
rectangle \(\prod_k(c_k,e_k)\) with real \(c_k<e_k\) is the union over
large \(j\) of the finite rectangles \(\prod_k(c_k,e_k-1/j]\), so by
part~(ii) of \cref{thm:p15-w02c-p15-002-generator-variants}
\(\mathcal L=\mathcal B(\mathbb R^d)\). Hence
\(\sigma(\mathcal S_n)=\mathcal B(B_n)\).

\emph{Step 5: gluing.} The measures \(E\mapsto\mu_{n+1}(E)\) and \(\mu_n\)
on \(\mathcal B(B_n)\) agree on the pi-system \(\mathcal S_n\), both being
\(m\) there, and \(B_n\in\mathcal S_n\) has finite measure; by
\cref{thm:p15-w02c-p15-002-sigma-finite-uniqueness} with \(A_k=B_n\), they
agree on \(\mathcal B(B_n)\). For Borel \(E\), put
\(\mu(E)=\lim_n\mu_n(E\cap B_n)\); the limit exists in \([0,\infty]\)
because \(\mu_{n+1}(E\cap B_{n+1})\geq\mu_{n+1}(E\cap B_n)=\mu_n(E\cap B_n)\).
Clearly \(\mu(\varnothing)=0\). If \(E_1,E_2,\ldots\) are disjoint with
union \(E\), then
\(\mu_n(E\cap B_n)=\sum_k\mu_n(E_k\cap B_n)\), each term nondecreasing in
\(n\). The limit of the left side is at least
\(\sum_{k\leq K}\lim_n\mu_n(E_k\cap B_n)\) for every \(K\), and each
\(\sum_k\mu_n(E_k\cap B_n)\) is at most \(\sum_k\mu(E_k)\); hence
\(\mu(E)=\sum_k\mu(E_k)\), and \(\mu\) is a measure. A finite rectangle
\(A\) lies in \(B_n\) for all large \(n\), and then
\(\mu_n(A\cap B_n)=m(A)\); so \(\mu(A)=\Delta_AF\). A bounded Borel set
lies in some \(B_n\), which has finite measure \(m(B_n)\).

\emph{Uniqueness in (i).} The class \(\mathcal P\) of finite rectangles and
\(\varnothing\) is a pi-system generating \(\mathcal B(\mathbb R^d)\)
(Step~4), and \(B_n\in\mathcal P\) increase to \(\mathbb R^d\) with finite
measure. Two measures with the stated rectangle values agree on
\(\mathcal B(\mathbb R^d)\) by
\cref{thm:p15-w02c-p15-002-sigma-finite-uniqueness}.

(ii) Fix \(x\). The rectangles \((x-n\mathbf1,x]\) increase to
\((-\infty,x]\), so by part~(iii) of \cref{thm:p15-measure-laws},
\(\mu((-\infty,x])=\lim_n\Delta_{(x-n\mathbf1,x]}F\). The vertex \(x\)
has sign \(+1\). Every other vertex \(v\) has some coordinate
\(v_i=x_i-n\) and all coordinates at most those of \(x\), so by (i)
\(0\leq F(v)\leq F(x[i\gets x_i-n])\), which tends to \(0\) by
(iii\(^*\)). There are \(2^d-1\) such vertices, so
\(\mu((-\infty,x])=F(x)\). Then \((-\infty,n\mathbf1]\) increase to
\(\mathbb R^d\) and \(\mu(\mathbb R^d)=\lim_nF(n\mathbf1)=1\) by
(iii\(^*\)).

(iii) If \(F\) is the distribution function of a probability measure, it
satisfies (i), (ii), (iii\(^*\)) and (iv) by
\cref{thm:p15-w02c-p15-002-rectangle-inclusion-exclusion}, and the measure
gives each finite rectangle \(A\) the mass \(\Delta_AF\), so it is unique
by (i). Conversely (i) and (ii) give a probability measure with
distribution function \(F\).
\end{proof}
\end{annexitem}

\begin{annexitem}{annex:p15-w02c-p15-002-product-difference}{thm:p15-w02c-p15-002-product-difference}{Differences of a product}
\begin{proof}
For \(G(x)=\prod_iF_i(x_i)\),
\((D^k_{s,t}G)(x)=(F_k(t)-F_k(s))\prod_{i\neq k}F_i(x_i)\): only the
\(k\)-th factor changes. Applying the operators of \((\dagger)\) one
coordinate at a time replaces each factor \(F_k(x_k)\) by
\(F_k(b_k)-F_k(a_k)\), so
\(\Delta_AF=\prod_k(F_k(b_k)-F_k(a_k))\). If each \(F_k\) is nondecreasing,
each factor is nonnegative, which is (iv). If each \(F_k\) is right
continuous and \(y\downarrow x\), each \(F_k(y_k)\to F_k(x_k)\), and a
product of finitely many convergent real sequences converges to the
product of the limits; this is (ii). If each \(F_k\) is a cumulative
distribution function, it is nondecreasing with values in \([0,1]\),
\(F_k(t)\to0\) as \(t\to-\infty\) and \(F_k(t)\to1\) as \(t\to\infty\)
(\cref{thm:p15-cdf-laws}). Then \(F\) takes values in \([0,1]\) and is
nondecreasing, a product of nonnegative nondecreasing factors.
\(F(x[i\gets t])=F_i(t)\prod_{k\neq i}F_k(x_k)\to0\) as \(t\to-\infty\),
and \(F(x_n)=\prod_kF_k(x_{n,k})\to1\) when every coordinate tends to
\(\infty\). This is (iii\(^*\)).
\end{proof}
\end{annexitem}

\begin{annexitem}{annex:p15-w02c-p15-002-indicator-membership}{thm:p15-w02c-p15-002-indicator-membership}{Indicators and membership}
\begin{proof}
For a Borel set \(B\), the inverse image \(\mathbf1_A^{-1}(B)\) is
determined by which of the values \(0,1\) lie in \(B\): it is \(\Omega\)
if both do, \(A\) if only \(1\) does, \(A^c\) if only \(0\) does, and
\(\varnothing\) if neither does. The Borel sets \(\mathbb R\), \(\{1\}\),
\(\{0\}\) and \(\varnothing\) realize all four cases, so
\(\sigma(\mathbf1_A)=\{\varnothing,A,A^c,\Omega\}\).
By \cref{def:p15-measurable-map}, \(\mathbf1_A\) is
\(\mathcal F\)-measurable exactly when
\(\sigma(\mathbf1_A)\subseteq\mathcal F\). If \(A\in\mathcal F\), the four
sets lie in \(\mathcal F\), because a sigma-algebra contains
\(\varnothing\) and \(\Omega\) and is closed under complements.
If \(\sigma(\mathbf1_A)\subseteq\mathcal F\), then in particular
\(A\in\mathcal F\).
\end{proof}
\end{annexitem}

\begin{annexitem}{annex:p15-w02c-p15-002-inverse-sampling}{thm:p15-w02c-p15-002-inverse-sampling}{The inverse applied to any uniform variable}
\begin{proof}
Fix \(x\in\mathbb R\). For every \(\omega\in\Omega\) the value
\(u=U(\omega)\) lies in \((0,1)\), and \cref{thm:p15-cdf-existence} gives
\(Q_F(u)\leq x\) if and only if \(u\leq F(x)\). Applied at each outcome,
this is the set identity \(\{X\leq x\}=\{U\leq F(x)\}\).

\emph{Measurability.} The set \(\{U\leq F(x)\}\) is the inverse image under
\(U\) of the Borel ray \((-\infty,F(x)]\), so it lies in \(\mathcal F\)
because \(U\) is a random variable. Thus \(X^{-1}((-\infty,x])\in\mathcal F\)
for every real \(x\). These rays generate \(\mathcal B(\mathbb R)\)
(\cref{thm:p15-borel-generators}), so \(X\) is a real random variable by
\cref{thm:p15-measurability-criterion}.

\emph{The distribution function.} Three cases cover every \(x\), since
\(F(x)\in[0,1]\). If \(F(x)=0\), the event \(\{U\leq0\}\) is empty because
\(U>0\) everywhere, so \(P(X\leq x)=0=F(x)\). If \(0<F(x)<1\), the
hypothesis on \(U\) with \(u=F(x)\) gives \(P(X\leq x)=F(x)\). If
\(F(x)=1\), the event \(\{U\leq1\}\) is \(\Omega\) because \(U<1\)
everywhere, so \(P(X\leq x)=1=F(x)\).

\emph{The law.} By \cref{thm:p15-cdf-existence}, \(F\) is also the
cumulative distribution function of \(Q_F\) on the constructed uniform
space, whose law is \(\mu_F\). Two real random variables with the same
cumulative distribution function have the same law
(\cref{thm:p15-cdf-laws}), so the law of \(X\) is \(\mu_F\).
\end{proof}
\end{annexitem}

\begin{annexitem}{annex:p15-w02c-p15-002-equal-in-distribution}{thm:p15-w02c-p15-002-equal-in-distribution}{Equality in distribution through the uniqueness of Stieltjes measures}
\begin{proof}
\emph{Item (1).} By \cref{thm:p15-cdf-laws}, \(F_X\) is real-valued,
nondecreasing and right continuous, so it is a Stieltjes measure function
in the sense of \cref{def:p15-stieltjes-intervals}. Because \(F_X\) is
nondecreasing, its infimum over \(\mathbb R\) is its limit at
\(-\infty\) and its supremum is its limit at \(\infty\); by
\cref{thm:p15-cdf-laws} these limits are \(0\) and \(1\), so
\(F_X(-\infty)=0\) and \(F_X(\infty)=1\).

The law \(\mathbb P_X(B)=\mathbb P(X\in B)\) is a probability measure on
\(\mathcal B(\mathbb R)\) (\cref{def:p15-pushforward-law}). Its values on
the members of \(\mathcal S_{\mathbb R}\) are computed in four cases. For
real \(a<b\), \(\mathbb P_X((a,b])=\mathbb P(a<X\leq b)=F_X(b)-F_X(a)\) by
\cref{thm:p15-cdf-laws}. For \(a=-\infty\) and real \(b\),
\(I(-\infty,b)=(-\infty,b]\) and
\(\mathbb P_X((-\infty,b])=F_X(b)=F_X(b)-F_X(-\infty)\). For real \(a\)
and \(b=\infty\), \(I(a,\infty)=\mathbb R\setminus(-\infty,a]\), so
\(\mathbb P_X(I(a,\infty))=1-F_X(a)=F_X(\infty)-F_X(a)\). For
\(a=-\infty\) and \(b=\infty\),
\(\mathbb P_X(\mathbb R)=1=F_X(\infty)-F_X(-\infty)\). Thus
\(\mathbb P_X\) is a Borel measure on \(\mathbb R\) with
\(\mathbb P_X(I(a,b))=F_X(b)-F_X(a)\) for all allowed endpoints. By
\cref{thm:p15-stieltjes-construction} there is exactly one Borel measure
with this property, namely \(\mu_{F_X}\); hence \(\mathbb P_X=\mu_{F_X}\).

\emph{Item (2).} Suppose \(X\stackrel{d}{=}Y\), that is,
\(\mathbb P_X=\mathbb P'_Y\) on \(\mathcal B(\mathbb R)\). Each ray
\((-\infty,x]\) is a Borel set, so
\(F_X(x)=\mathbb P_X((-\infty,x])=\mathbb P'_Y((-\infty,x])=F_Y(x)\).
Conversely, suppose \(F_X=F_Y=F\). Item~(1), applied to \(X\) and to
\(Y\) (its proof uses only that each is a real random variable on some
probability space), gives \(\mathbb P_X=\mu_F\) and \(\mathbb P'_Y=\mu_F\).
Hence \(\mathbb P_X=\mathbb P'_Y\), which is \(X\stackrel{d}{=}Y\).
\end{proof}
\end{annexitem}

\begin{annexitem}{annex:p15-w02c-p15-002-dense-jumps}{thm:p15-w02c-p15-002-dense-jumps}{The law behind Example 1.2.9}
\begin{proof}
For a Borel set \(B\) write \(I_B=\{i\geq1:q_i\in B\}\), so that
\(\mu(B)=\sum_{i\in I_B}\alpha_i\), a sum of nonnegative terms over a
countable index set, with value in \([0,1]\) because the full sum is one.

(i) \emph{Measure.} The empty set gives the empty sum, so
\(\mu(\varnothing)=0\). Let \(B_1,B_2,\ldots\) be pairwise disjoint Borel sets
with union \(B\). The index sets \(I_{B_n}\) are pairwise disjoint, since
\(q_i\) lies in at most one \(B_n\), and their union is \(I_B\). The grouping
statement of \cref{thm:p15-countable-atomic}, applied to the countable set
\(I_B\) and its partition into the sets \(I_{B_n}\), gives
\[
 \mu(B)=\sum_{i\in I_B}\alpha_i
       =\sum_{n\geq1}\sum_{i\in I_{B_n}}\alpha_i
       =\sum_{n\geq1}\mu(B_n).
\]
So \(\mu\) is a measure, and \(\mu(\mathbb R)=\sum_{i\geq1}\alpha_i=1\).
Every \(q_i\) is rational, so \(I_{\mathbb R\setminus\mathbb Q}\) is empty and
\(\mu(\mathbb R\setminus\mathbb Q)=0\); as \(\mathbb Q\) is countable, \(\mu\)
is discrete with carrier \(\mathbb Q\).

\emph{Distribution function.} Fix \(x\). The term
\(\alpha_i\mathbf1_{[q_i,\infty)}(x)\) equals \(\alpha_i\) when
\(q_i\leq x\) and \(0\) otherwise. Grouping the indices into
\(I_{(-\infty,x]}\) and its complement (again by
\cref{thm:p15-countable-atomic}) gives
\[
 F(x)=\sum_{i\in I_{(-\infty,x]}}\alpha_i+0=\mu((-\infty,x]).
\]
Let \(X\) be the identity map of \(\mathbb R\), a real random variable on
\((\mathbb R,\mathcal B(\mathbb R),\mu)\) because \(X^{-1}(B)=B\). Its law is
\(\mu\) and its cumulative distribution function is \(F_X=F\). Hence
\cref{thm:p15-cdf-laws} shows that \(F\) is nondecreasing and right
continuous with limits \(0\) at \(-\infty\) and \(1\) at \(\infty\).

(ii) By \cref{thm:p15-cdf-laws}, \(F(x)-F(x-)=p_X(x)=\mu(\{x\})\) for every
\(x\). If \(x=q_i\), then \(I_{\{x\}}=\{i\}\), because the enumeration lists
each rational exactly once, so \(\mu(\{x\})=\alpha_i\). If \(x\) is
irrational, \(I_{\{x\}}\) is empty and \(\mu(\{x\})=0\). The same theorem
identifies the discontinuities of \(F_X\) with the points of positive mass;
since every \(\alpha_i\) is positive, these are exactly the rationals. Let
\(a<b\). Choose a rational \(r_1\in(a,b)\), and, given rationals
\(r_1>\cdots>r_k\) in \((a,b)\), a rational \(r_{k+1}\in(a,r_k)\). The
\(r_k\) are distinct rationals in \((a,b)\), so the interval contains
infinitely many discontinuities of \(F\).

(iii) Let \(x<y\) and choose a rational \(r\in(x,y)\); then \(r=q_i\) for some
\(i\). The sets \((-\infty,x]\) and \((x,y]\) are disjoint with union
\((-\infty,y]\), so by (i) and additivity of \(\mu\),
\[
 F(y)-F(x)=\mu((x,y])\geq\mu(\{q_i\})=\alpha_i>0,
\]
the inequality being monotonicity of \(\mu\) under
\(\{q_i\}\subseteq(x,y]\).
\end{proof}
\end{annexitem}

\begin{annexitem}{annex:p15-w02c-p15-002-order-borel}{thm:p15-w02c-p15-002-order-borel}{Borel sets of the extended line}
\begin{proof}
Let \(\mathfrak B\) consist of the intervals \([-\infty,a)\), \((a,b)\)
and \((b,\infty]\) of \(\mathbf R^*\) with real endpoints, and let
\(\mathfrak B_{\mathbb Q}\) be its members with rational endpoints.

\emph{The generated topology.} The intersection of two members of
\(\mathfrak B\) is empty or again a member: two left rays meet in the
shorter one, two right rays in the shorter one, and in every other case
the intersection is an interval \((c,d)\) whose endpoints are the larger
of the left endpoints and the smaller of the right endpoints, or is empty
when these are out of order. Also
\(\mathbf R^*=[-\infty,1)\cup(0,\infty]\). Let \(\tau\) be the collection
of all unions of members of \(\mathfrak B\), the empty union included. It
is closed under arbitrary unions; the intersection of two unions is the
union of the pairwise intersections, each empty or in \(\mathfrak B\); and
it contains \(\mathbf R^*\). So \(\tau\) is a topology containing
\(\mathfrak B\), and every topology containing \(\mathfrak B\) contains all
its unions. Hence \(\tau\) is the topology generated by \(\mathfrak B\),
and \(\mathcal R^*=\sigma(\tau)\).

\emph{Rational refinement.} Let \(W\in\mathfrak B\) and \(x\in W\). If
\(W=[-\infty,a)\), choose a rational \(r\) with \(x<r<a\) (any rational
below \(a\) when \(x=-\infty\)); then \(x\in[-\infty,r)\subseteq W\). If
\(W=(b,\infty]\), choose a rational \(q\) with \(b<q<x\) (any rational
above \(b\) when \(x=\infty\)); then \(x\in(q,\infty]\subseteq W\). If
\(W=(a,b)\), then \(x\) is real; choose rationals \(a<q<x<r<b\); then
\(x\in(q,r)\subseteq W\). Now let \(U\in\tau\). Each \(x\in U\) lies in
some \(W\in\mathfrak B\) with \(W\subseteq U\), and hence in some
\(W'\in\mathfrak B_{\mathbb Q}\) with \(W'\subseteq U\). Therefore \(U\) is
the union of the members of \(\mathfrak B_{\mathbb Q}\) contained in it.
Since \(\mathfrak B_{\mathbb Q}\) is indexed by rationals or pairs of
rationals, it is countable, and \(U\) is a countable union of its members.

\emph{Comparison.} Let \(\mathcal G=\sigma(\{[-\infty,a]:a\in\mathbb R\})\).
For rationals \(q<r\),
\[
 [-\infty,q)=\bigcup_{n\geq1}[-\infty,q-1/n],\qquad
 (r,\infty]=\mathbf R^*\setminus[-\infty,r],
\]
\[
 (q,r)=[-\infty,r)\setminus[-\infty,q],
\]
so \(\mathfrak B_{\mathbb Q}\subseteq\mathcal G\). By the refinement step
every set in \(\tau\) is a countable union of members of
\(\mathfrak B_{\mathbb Q}\), so \(\tau\subseteq\mathcal G\) and
\(\mathcal R^*=\sigma(\tau)\subseteq\mathcal G\). Conversely
\([-\infty,a]=\mathbf R^*\setminus(a,\infty]\) is the complement of an open
set, so it lies in \(\mathcal R^*\), and \(\mathcal G\subseteq\mathcal R^*\).

The two sigma-algebras on the same ordered set are equal, so
\(\mathcal F/\mathcal R^*\)-measurability and measurability in
\cref{def:p15-extended-random-variables} are the same condition, and
every statement proved for that measurable space applies unchanged.
\end{proof}
\end{annexitem}

\begin{annexitem}{annex:p15-w02c-p15-002-domain-extension}{thm:p15-w02c-p15-002-domain-extension}{Random variables on an event}
\begin{proof}
\emph{The trace collection.} \(D\in\mathcal F_D\). If \(A\in\mathcal F_D\),
then \(D\setminus A\in\mathcal F\) and \(D\setminus A\subseteq D\). A
countable union of members of \(\mathcal F_D\) lies in \(\mathcal F\) and
in \(D\). So \(\mathcal F_D\) is a sigma-algebra on \(D\). Each
\(A\in\mathcal F_D\) equals \(D\cap A\); conversely, for \(A\in\mathcal F\)
the set \(D\cap A\) lies in \(\mathcal F\) and in \(D\). Hence
\(\mathcal F_D=\{D\cap A:A\in\mathcal F\}\).

(i)\(\Leftrightarrow\)(ii). Every inverse image \(X^{-1}(B)\) is a subset
of \(D\), so it lies in \(\mathcal F\) exactly when it lies in
\(\mathcal F_D\).

(i)\(\Rightarrow\)(iii). For \(B\in\mathcal R^*\),
\[
 (X^c)^{-1}(B)=X^{-1}(B)\cup E_B,\qquad
 E_B=\begin{cases}\Omega\setminus D,&c\in B,\\ \varnothing,&c\notin B,\end{cases}
\]
and both sets lie in \(\mathcal F\). So \(X^c\) is
\(\mathcal F/\mathcal R^*\)-measurable, which is the definition of an
extended real random variable by \cref{thm:p15-w02c-p15-002-order-borel}.

(iii)\(\Rightarrow\)(i). For \(B\in\mathcal R^*\),
\(X^{-1}(B)=(X^c)^{-1}(B)\cap D\in\mathcal F\).
\end{proof}
\end{annexitem}

\begin{annexitem}{annex:p15-w02c-p15-002-convergence-set}{thm:p15-w02c-p15-002-convergence-set}{The convergence set is an event, and the two forms of almost sure convergence}
\begin{proof}
Fix a real sequence \((x_n)\) and put \(i_k=\inf_{n\geq k}x_n\) and
\(s_k=\sup_{n\geq k}x_n\) in \(\mathbf R^*\). The sequence \(i_k\)
increases and \(s_k\) decreases in \(k\), and
\(L=\sup_ki_k\), \(U=\inf_ks_k\). A real sequence converges to a real
\(c\) when for every \(\varepsilon>0\) eventually \(|x_n-c|<\varepsilon\);
it converges to \(\infty\) when for every real \(M\) eventually
\(x_n>M\); and to \(-\infty\) when for every real \(M\) eventually
\(x_n<M\).

(i) For any \(k,j\), with \(n=\max(k,j)\),
\(i_k\leq x_n\leq s_j\). Taking the supremum over \(k\) gives
\(L\leq s_j\) for each \(j\), and the infimum over \(j\) gives
\(L\leq U\).

Suppose \(L=U=c\). If \(c\) is real and \(\varepsilon>0\), there are
\(k_1\) with \(i_{k_1}>c-\varepsilon\) (since \(\sup_ki_k=c\)) and \(k_2\)
with \(s_{k_2}<c+\varepsilon\) (since \(\inf_ks_k=c\)); for
\(n\geq\max(k_1,k_2)\),
\(c-\varepsilon<i_{k_1}\leq x_n\leq s_{k_2}<c+\varepsilon\). If
\(c=\infty\) and \(M\) is real, some \(i_k>M\), and then \(x_n>M\) for
\(n\geq k\). The case \(c=-\infty\) is the same argument applied to
\(U=-\infty\) and the \(s_k\).

Conversely suppose \(x_n\to c\). If \(c\) is real and \(\varepsilon>0\),
choose \(N\) with \(|x_n-c|<\varepsilon\) for \(n\geq N\); then
\(c-\varepsilon\leq i_N\leq L\leq U\leq s_N\leq c+\varepsilon\). As
\(\varepsilon\) is arbitrary, \(L=U=c\). If \(c=\infty\), for each real
\(M\) some \(i_N\geq M\), so \(L\geq M\) for every \(M\); thus \(L=\infty\)
and \(U\geq L\) gives \(U=\infty\). The case \(c=-\infty\) is symmetric,
with \(U\leq M\) for every \(M\).

(ii) By \cref{thm:p15-measurable-limits} and
\cref{thm:p15-w02c-p15-002-order-borel}, \(L\) and \(U\) are extended
real random variables, so the sets \(\{L\leq a\}\) and \(\{U\leq a\}\)
are events for real \(a\). Hence
\(\{L<q\}=\bigcup_n\{L\leq q-1/n\}\) and
\(\{U>r\}=\Omega\setminus\{U\leq r\}\) are events. If \(L(\omega)<U(\omega)\),
there is a rational \(q\) strictly between them (if both are real, by
density; if \(L=-\infty\), any rational below \(U\); if \(U=\infty\), any
rational above \(L\)), and then a rational \(r\) with \(q<r<U(\omega)\) by
the same reasoning. Conversely \(L<q<r<U\) implies \(L<U\). This proves
the displayed union, a countable union of events. Since \(L\leq U\)
everywhere by (i), \(\{L=U\}\) is the complement of \(\{L<U\}\).

By (i), the limit exists in \(\mathbb R\) exactly when \(L=U\) and the
common value is real, that is, \(L=U\) and \(-m<U<m\) for some integer
\(m\geq1\). Each \(\{-m<U<m\}=\{U<m\}\setminus\{U\leq-m\}\), with
\(\{U<m\}=\bigcup_k\{U\leq m-1/k\}\), is an event, so \(\Omega_o\) is an
event. On \(\{L=U\}\) the limit is the common value, which is \(U=X_\infty\),
again by (i).

(iii) \(X_\infty=U\) is an extended real random variable by
\cref{thm:p15-measurable-limits}. For \(B\in\mathcal R^*\), the inverse
image of \(B\) under the restriction of \(X_\infty\) to \(D=\{L=U\}\) is
\(X_\infty^{-1}(B)\cap D\), an intersection of two events, using
\cref{thm:p15-w02c-p15-002-order-borel} for \(X_\infty^{-1}(B)\in\mathcal F\).

(iv) The indicator \(\mathbf1_{\Omega_o}\) is a real random variable by
\cref{thm:p15-w02c-p15-002-indicator-membership}, so each
\(Y_n=X_n\mathbf1_{\Omega_o}\) is one, as a product of real random
variables (\cref{thm:p15-measurable-operations}). For
\(\omega\in\Omega_o\), \(Y_n(\omega)=X_n(\omega)\) converges to the real
number \(X_\infty(\omega)=Y(\omega)\). For \(\omega\notin\Omega_o\),
\(Y_n(\omega)=0=Y(\omega)\) for every \(n\). So the real random variables
\(Y_n\) converge at every outcome to the finite value \(Y\), and \(Y\) is a
real random variable by \cref{thm:p15-measurable-limits}. The event
\(\{X_n\to Y\}\) of that theorem contains \(\Omega_o\), because on
\(\Omega_o\) the sequence \(X_n\) converges to \(X_\infty=Y\). By
monotonicity \(\mathbb P(X_n\to Y)\geq\mathbb P(\Omega_o)=1\).

(v) At every \(\omega\in\{X_n\to X\}\) the limit of \(X_n(\omega)\)
exists in \(\mathbb R\) and equals \(X(\omega)\), so
\(\{X_n\to X\}\subseteq\Omega_o\), and \(X(\omega)=X_\infty(\omega)\)
there by (i). Monotonicity gives \(\mathbb P(\Omega_o)=1\). If also
\(X_n\to X'\) almost surely, the event \(\{X_n\to X\}\cap\{X_n\to X'\}\)
has probability one, because its complement is a union of two null
events (\cref{thm:p15-measure-laws}), and on it
\(X=X_\infty=X'\).
\end{proof}
\end{annexitem}

\begin{annexitem}{annex:p15-w02c-p15-002-sigma-finite-uniqueness}{thm:p15-w02c-p15-002-sigma-finite-uniqueness}{Durrett's Theorem A.1.5: finite exhaustion in a pi-system}
\begin{proof}
All sets below lie in \(\sigma(\mathcal P)\), which is contained in
\(\mathcal F_1\cap\mathcal F_2\), so both measures are defined on them.

\emph{Step 1.} Fix \(A\in\mathcal P\) with \(\nu_1(A)<\infty\); then
\(\nu_2(A)=\nu_1(A)<\infty\). Let
\[
 \mathcal L_A=\{B\in\sigma(\mathcal P):\nu_1(A\cap B)=\nu_2(A\cap B)\}.
\]
This is a lambda-system on \(\Omega\) (\cref{def:p15-pi-lambda}).
First, \(\Omega\in\mathcal L_A\), since \(A\cap\Omega=A\in\mathcal P\).
Next, let \(C\subseteq B\) be members of \(\mathcal L_A\). Then
\(A\cap C\subseteq A\cap B\),
\(A\cap(B\setminus C)=(A\cap B)\setminus(A\cap C)\), and
\(\nu_i(A\cap B)\leq\nu_i(A)<\infty\) by monotonicity. Part~(i) of
\cref{thm:p15-measure-laws} gives
\[
 \nu_i(A\cap(B\setminus C))=\nu_i(A\cap B)-\nu_i(A\cap C)\qquad(i=1,2),
\]
and the right-hand sides agree, so \(B\setminus C\in\mathcal L_A\).
Finally, if \(B_n\in\mathcal L_A\) increase to \(B\), then
\(A\cap B_n\) increase to \(A\cap B\), and part~(iii) of
\cref{thm:p15-measure-laws} gives
\(\nu_1(A\cap B)=\lim_n\nu_1(A\cap B_n)=\lim_n\nu_2(A\cap B_n)=\nu_2(A\cap B)\).

The pi-system \(\mathcal P\) is contained in \(\mathcal L_A\): for
\(B\in\mathcal P\), \(A\cap B\in\mathcal P\), where the measures agree.
The first statement of \cref{thm:p15-measure-uniqueness} now gives
\(\sigma(\mathcal P)\subseteq\mathcal L_A\).

\emph{Step 2.} Let \(B\in\sigma(\mathcal P)\). Each \(A_n\) satisfies the
hypothesis of Step~1, so \(\nu_1(A_n\cap B)=\nu_2(A_n\cap B)\) for all
\(n\). Since \(A_n\) increase to \(\Omega\), the sets \(A_n\cap B\)
increase to \(B\), and part~(iii) of \cref{thm:p15-measure-laws} gives
\(\nu_1(B)=\lim_n\nu_1(A_n\cap B)=\lim_n\nu_2(A_n\cap B)=\nu_2(B)\).
\end{proof}
\end{annexitem}

\begin{annexitem}{annex:p15-w02c-p15-002-borel-paradox}{thm:p15-w02c-p15-002-borel-paradox}{The conditional law of the angle along the circle}
\begin{proof}
(i) Let \(\mathcal G=\sigma(\Theta')\) and \(K(\omega,A)=\mathbb P(\Psi\in A)\).
For fixed \(A\), \(K(\cdot,A)\) is constant, so each of its inverse images
is \(\varnothing\) or \(\Omega\), and it is \(\mathcal G\)-measurable.
For fixed \(\omega\), \(K(\omega,\cdot)\) is the law of \(\Psi\), a
probability measure on \((\mathbb R,\mathcal B(\mathbb R))\) by
\cref{thm:p15-pushforward-laws}. So \(K\) is a probability kernel from
\((\Omega,\mathcal G)\) to \((\mathbb R,\mathcal B(\mathbb R))\)
(\cref{def:p15-probability-kernel}). Every \(G\in\mathcal G\) has the
form \(\{\Theta'\in C\}\) with \(C\) Borel (\cref{def:p15-measurable-map}).
Since \(\Theta'\) and \(\Psi\) are independent
(\cref{thm:p15-conditioning-coordinate-examples},
\cref{def:p15-independent-random-elements}),
\[
 \int_G\mathbf1_{\{\Psi\in A\}}\,dP=\mathbb P(\Psi\in A,\ \Theta'\in C)
 =\mathbb P(\Psi\in A)\,\mathbb P(\Theta'\in C)=\int_GK(\cdot,A)\,dP .
\]
The function \(K(\cdot,A)\) is bounded and \(\mathcal G\)-measurable, so it
is a version of \(P(\Psi\in A\mid\mathcal G)\)
(\cref{def:p15-general-conditional-expectation}). This holds for every
Borel \(A\), so \(K\) is a regular conditional law of \(\Psi\) given
\(\mathcal G\) (\cref{def:p15-regular-conditional-law}). The same argument
with \(\Theta,\Phi\), which are independent by the same theorem, gives the
statement for \(\Phi\) given \(\sigma(\Theta)\).

(ii) \(\Psi\) is a real random variable. By the last statement of
\cref{thm:p15-real-conditional-law}, \(K'\) and \(K\) agree as probability
measures outside one \(\mathcal G\)-measurable null set, simultaneously
for every Borel event. Outside that set \(K'(\omega,\cdot)\) is the law of
\(\Psi\), which by (iii) gives \([-\pi/6,\pi/6]\) probability \(1/4\). The
uniform law on \((-\pi,\pi]\), the Borel probability
\(B\mapsto\lambda(B\cap(-\pi,\pi])/(2\pi)\), gives that interval
\((\pi/3)/(2\pi)=1/6\). So \(K'(\omega,\cdot)\) is not uniform there.

(iii) By \cref{thm:p15-conditioning-coordinate-examples}, \(\Phi\) has
density \(\cos\phi/2\) on \((-\pi/2,\pi/2]\) and \(\Psi\) has density
\(|\cos\psi|/4\) on \((-\pi,\pi]\). These densities are continuous on the
closed intervals used below, where the Lebesgue integral of a continuous
function equals its Riemann integral, and \(\sin\) is an antiderivative of
\(\cos\). Single points have probability zero under a law with a density.
Hence
\[
 \mathbb P(|\Phi|\leq\pi/6)=\int_{-\pi/6}^{\pi/6}\frac{\cos\phi}2\,d\phi
 =\frac{\sin(\pi/6)-\sin(-\pi/6)}2=\frac12,
\]
\[
 \mathbb P(|\Phi|\geq\pi/3)=1-\mathbb P(|\Phi|<\pi/3)
 =1-\frac{\sin(\pi/3)-\sin(-\pi/3)}2=1-\frac{\sqrt3}2,
\]
and, since \(\cos\psi\geq0\) on \([-\pi/6,\pi/6]\),
\[
 \mathbb P(|\Psi|\leq\pi/6)=\int_{-\pi/6}^{\pi/6}\frac{\cos\psi}4\,d\psi
 =\frac{2\sin(\pi/6)}4=\frac14 .
\]
The band \(|\phi|\leq\pi/6\) has length \(\pi/3\), and the two caps
\((-\pi/2,-\pi/3]\) and \([\pi/3,\pi/2]\) have total length \(\pi/3\);
each is one third of the length \(\pi\) of the latitude range. The
interval \([-\pi/6,\pi/6]\) has length \(\pi/3\), one sixth of \(2\pi\).
Finally \(1.72^2=2.9584<3\), so \(\sqrt3>1.72\) and
\(1-\sqrt3/2<1-0.86=0.14\).
\end{proof}
\end{annexitem}

\begin{annexitem}{annex:p15-w02c-p15-002-ae-additivity}{thm:p15-w02c-p15-002-ae-additivity}{Countable additivity outside a null set}
\begin{proof}
(i) Put \(B_m=\bigcup_{n\leq m}A_n\) and \(B=\bigcup_nA_n\). Because the
\(A_n\) are disjoint,
\(\mathbf1_{\{X\in B_m\}}=\sum_{n\leq m}\mathbf1_{\{X\in A_n\}}\). By the
linearity in \cref{thm:p15-conditional-basic}, the
\(\mathcal G\)-measurable function \(\sum_{n\leq m}g_{A_n}\) is a version
of \(P(X\in B_m\mid\mathcal G)\). Two versions agree outside a null set
(\cref{thm:p15-conditional-existence}), and the set where two
\(\mathcal G\)-measurable functions differ lies in \(\mathcal G\); so there
is a \(\mathcal G\)-measurable null set \(N_m\) off which
\(g_{B_m}=\sum_{n\leq m}g_{A_n}\).

The indicators \(\mathbf1_{\{X\in B_m\}}\) are nonnegative and increase to
\(\mathbf1_{\{X\in B\}}\) at every outcome. Part~(i) of
\cref{thm:p15-conditional-limits}, applied to the versions \(g_{B_m}\) and
\(g_B\), gives \(g_{B_m}\uparrow g_B\) outside a null set. The set \(N_0\)
of outcomes where \((g_{B_m}(\omega))_m\) fails to increase to
\(g_B(\omega)\) is defined from countably many \(\mathcal G\)-measurable
functions by countable operations, so \(N_0\in\mathcal G\), and it is null.

Let \(N=N_0\cup\bigcup_mN_m\), a \(\mathcal G\)-measurable null set by
part~(ii) of \cref{thm:p15-measure-laws}. For \(\omega\notin N\), the
partial sums \(\sum_{n\leq m}g_{A_n}(\omega)=g_{B_m}(\omega)\) increase to
\(g_B(\omega)\), so the series converges and its sum is
\(g_B(\omega)\).

(ii) Each sequence in the family has its null set from (i). A countable
union of \(\mathcal G\)-measurable null sets is \(\mathcal G\)-measurable
and null, by part~(ii) of \cref{thm:p15-measure-laws}.
\end{proof}
\end{annexitem}

\begin{annexitem}{annex:p15-w02c-p15-002-nice-rcd}{thm:p15-w02c-p15-002-nice-rcd}{Theorems 4.1.17 and 4.1.18 for nice spaces whose law is carried by a Borel set}
\begin{proof}
If \(\varphi(S)\) is Borel, \(E=\varphi(S)\) is a Borel subset of the image
with \(P(\varphi(X)\in E)=1\). In general let \(E\) be as in the
hypothesis.

\emph{A coded subspace.} Put \(S_E=\varphi^{-1}(E)\), a member of
\(\mathcal S\) because \(\varphi\) is measurable, and give it the trace
sigma-algebra \(\mathcal S_E=\{A\cap S_E:A\in\mathcal S\}\). Since
\(E\subseteq\varphi(S)\) and \(\varphi\) is one-to-one, \(\varphi\)
restricts to a bijection \(\psi:S_E\to E\). It is measurable into
\((E,\mathcal B(E))\): \(\psi^{-1}(E\cap C)=S_E\cap\varphi^{-1}(C)\) for
Borel \(C\). Its inverse is measurable: for \(A\in\mathcal S\), the
measurability of \(\varphi^{-1}\) on \(\varphi(S)\) gives a Borel \(C\)
with \(\varphi(A)=\varphi(S)\cap C\), and then
\(\psi(A\cap S_E)=\varphi(A)\cap E=C\cap E\in\mathcal B(E)\). So \(\psi\) is
a Borel real coding of \((S_E,\mathcal S_E)\), which is a standard Borel
space in the sense of \cref{def:p15-measurable-real-coding}.

\emph{A modified observation.} The event
\(\Omega_E=\{X\in S_E\}=\{\varphi(X)\in E\}\) has probability one, so it is
nonempty; fix \(\omega_0\in\Omega_E\) and \(s_0=X(\omega_0)\in S_E\). Let
\(X'=X\) on \(\Omega_E\) and \(X'=s_0\) off \(\Omega_E\). For
\(A\in\mathcal S\), the set \(\{X'\in A\cap S_E\}\) is
\(\{X\in A\}\cap\Omega_E\) if \(s_0\notin A\), and
\((\{X\in A\}\cap\Omega_E)\cup(\Omega\setminus\Omega_E)\) if \(s_0\in A\);
in both cases an event. So \(X'\) is measurable into
\((S_E,\mathcal S_E)\). Moreover \(\{X'\in A\cap S_E\}\) and
\(\{X\in A\}\) differ only inside the null event
\(\Omega\setminus\Omega_E\), so their indicators agree almost surely, and
by \cref{def:p15-general-conditional-expectation} every version of
\(P(X'\in A\cap S_E\mid\mathcal G)\) is a version of
\(P(X\in A\mid\mathcal G)\): the defining integrals over members of
\(\mathcal G\) are the same.

(i) By \cref{thm:p15-standard-conditional-law} applied to \(X'\), there is
a probability kernel \(K'\) from \((\Omega,\mathcal G)\) to
\((S_E,\mathcal S_E)\) with \(K'(\cdot,A')\) a version of
\(P(X'\in A'\mid\mathcal G)\) for every \(A'\in\mathcal S_E\). Define
\(K(\omega,A)=K'(\omega,A\cap S_E)\) for \(A\in\mathcal S\). For fixed
\(\omega\), \(A\mapsto A\cap S_E\) preserves disjoint countable unions and
sends \(S\) to \(S_E\), so \(K(\omega,\cdot)\) is a probability measure on
\(\mathcal S\); for fixed \(A\), \(K(\cdot,A)\) is \(\mathcal G\)-measurable
and, by the previous paragraph, a version of \(P(X\in A\mid\mathcal G)\).
So \(K\) is a regular conditional law, and it satisfies (i) and (ii) of the
source's definition, (ii) at every \(\omega\).

(ii) By \cref{thm:p15-value-conditional-law} applied to \(X'\) and \(Y\),
there is a probability kernel \(K'\) from \((T,\mathcal T)\) to
\((S_E,\mathcal S_E)\) with \(K'(Y,A')\) a version of
\(P(X'\in A'\mid\sigma(Y))\). Put \(K(y,A)=K'(y,A\cap S_E)\). As in (i),
\(K\) is a probability kernel from \(T\) to \(S\), and \(K(Y,A)\) is a
version of \(P(X\in A\mid\sigma(Y))\) for each \(A\in\mathcal S\). For
\((T,\mathcal T)=(S,\mathcal S)\) these are properties (i) and (ii) of
Theorem~4.1.18 with \(\mu=K\).
\end{proof}
\end{annexitem}

\begin{annexitem}{annex:p15-w02c-p15-002-nice-rcd-failure}{thm:p15-w02c-p15-002-nice-rcd-failure}{Without a Borel carrier both theorems can fail}
\begin{proof}
Let \(H_U\), \(H_U^c=(0,1)\setminus H_U\), \(\mathcal F_H\) and \(P_H\) be as
in \cref{thm:p15-w02c-p15-008-thick-set} and
\cref{def:p15-w02c-p15-008-model}, and let
\((\Omega,\mathcal F,P)=((0,1),\mathcal F_H,P_H)\), a probability space by
part~(ii) of \cref{thm:p15-w02c-p15-008-model}. Members of \(\mathcal F_U\)
are the Borel subsets of \((0,1)\), and every member of \(\mathcal F_H\)
has the form \((B_1\cap H_U)\cup(B_2\cap H_U^c)\) with
\(B_1,B_2\in\mathcal F_U\).

\emph{The space.} Let \(S=(0,1)\times\{0,1\}\) and
\(\mathcal S=\{(E_0\times\{0\})\cup(E_1\times\{1\}):E_0\in\mathcal F_H,\
E_1\in\mathcal F_U\}\). It is a sigma-algebra: complements and countable
unions act on the two copies separately, and \(\mathcal F_H\) and
\(\mathcal F_U\) are sigma-algebras on \((0,1)\) (part~(i) of
\cref{thm:p15-w02c-p15-008-model}). Define \(X(\omega)=(\omega,0)\) and
\(Y(\omega)=(\omega,1)\). The inverse image of
\((E_0\times\{0\})\cup(E_1\times\{1\})\) is \(E_0\) under \(X\) and \(E_1\)
under \(Y\). Since \(\mathcal F_U\subseteq\mathcal F_H\), both maps are
measurable, and \(\sigma(Y)=\mathcal F_U\).

\emph{\(S\) is nice.} Let \(\varphi(\omega,0)=\omega\) for
\(\omega\in H_U^c\), \(\varphi(\omega,0)=\omega+2\) for \(\omega\in H_U\),
and \(\varphi(\omega,1)=\omega+4\). The three pieces take values in the
disjoint intervals \((0,1)\), \((2,3)\) and \((4,5)\), and \(\varphi\) is
one-to-one on each, so it is one-to-one. For Borel \(C\),
\(\varphi^{-1}(C)=(E_0\times\{0\})\cup(E_1\times\{1\})\) with
\[
 E_0=\bigl((C-2)\cap(0,1)\cap H_U\bigr)\cup\bigl(C\cap(0,1)\cap H_U^c\bigr)
 \in\mathcal F_H,
 \qquad E_1=(C-4)\cap(0,1)\in\mathcal F_U ;
\]
so \(\varphi\) is measurable. Let \(Z=\varphi(S)\); then
\(Z\cap(0,1)=H_U^c\), \(Z\cap(2,3)=H_U+2\) and \(Z\cap(4,5)=(4,5)\). For
\(E_0=(B_1\cap H_U)\cup(B_2\cap H_U^c)\) with \(B_1,B_2\in\mathcal F_U\) and
for \(E_1\in\mathcal F_U\),
\[
 \varphi\bigl((E_0\times\{0\})\cup(E_1\times\{1\})\bigr)
 =Z\cap\bigl(B_2\cup(B_1+2)\cup(E_1+4)\bigr),
\]
because \(Z\cap B_2=H_U^c\cap B_2\), \(Z\cap(B_1+2)=(H_U\cap B_1)+2\) and
\(Z\cap(E_1+4)=E_1+4\). The set \(B_2\cup(B_1+2)\cup(E_1+4)\) is Borel.
Since \(\varphi\) is a bijection onto \(Z\), these images are the inverse
images of the sets of \(\mathcal S\) under \(\varphi^{-1}\); so
\(\varphi^{-1}\) is measurable from \(Z\) with its trace sigma-algebra,
and \(\varphi\) is a nice map.

\emph{No regular conditional distribution.} Let
\(\mathcal G=\sigma(Y)=\mathcal F_U\), and suppose \(\mu\) is a regular
conditional distribution for \(X\) given \(\mathcal G\) in the sense of
\cref{def:p15-w02b-p15-008-rcd}. For \(A\in\mathcal F_H\) the set
\(A\times\{0\}\) belongs to \(\mathcal S\), and put
\(\tilde\mu(\omega,A)=\mu(\omega,A\times\{0\})\). Since
\(\{X\in A\times\{0\}\}=A\), the function \(\tilde\mu(\cdot,A)\) is a
version of \(P_H(A\mid\mathcal G)\), which is condition~(i) of that
definition for the identity map of \(((0,1),\mathcal F_H)\). In particular
\(\tilde\mu(\cdot,(0,1))\) and the constant \(1\) are both versions of
\(P_H((0,1)\mid\mathcal G)\), so they agree outside a null set
(\cref{thm:p15-conditional-existence}). Outside the union of that null set
and the null set where \(\mu(\omega,\cdot)\) is not a probability measure,
a null set by part~(ii) of \cref{thm:p15-measure-laws}, the set function
\(A\mapsto\tilde\mu(\omega,A)\) is countably additive, because
\(A\mapsto A\times\{0\}\) preserves disjoint countable unions, and has
total mass \(1\); so it is a probability measure on \(\mathcal F_H\). Thus
\(\tilde\mu\) is a regular conditional distribution for the identity map
of \(((0,1),\mathcal F_H)\) given \(\mathcal F_U\), which part~(i) of
\cref{thm:p15-w02c-p15-008-no-rcp} excludes. So no regular conditional
distribution for \(X\) given \(\sigma(Y)\) exists.

\emph{Theorem 4.1.18 fails.} If \(\mu':S\times\mathcal S\to[0,1]\) had
properties (i) and (ii) of Theorem~4.1.18, then
\(\mu(\omega,A)=\mu'(Y(\omega),A)\) would be a regular conditional
distribution for \(X\) given \(\sigma(Y)\) in the source's sense, which
was just excluded.

\emph{No Borel carrier.} If for some nice map \(\varphi'\) a Borel
\(E\subseteq\varphi'(S)\) had \(P(\varphi'(X)\in E)=1\), part~(i) of
\cref{thm:p15-w02c-p15-002-nice-rcd} would give a regular conditional law,
hence a regular conditional distribution. So no such \(E\) exists; taking
\(E=\varphi'(S)\) shows that \(\varphi'(S)\) is not Borel.
\end{proof}
\end{annexitem}
